\documentclass[french,11pt]{book}
\usepackage{teach}
\usepackage[french]{babel}
\usepackage{amsmath,amssymb}
\usepackage{eurosym}
\newcommand{\pourcent}{\%}

\begin{document}

\chapter*{Taux d'evolution, pourcentages et informations chiffrees}

\section*{Proportion et pourcentage}

Une proportion compare une partie a un tout. Si une partie a pour effectif $p$ dans un ensemble d'effectif total $T$, alors la proportion est
\[
\frac{p}{T}.
\]
Pour l'exprimer en pourcentage, on multiplie par $100$.

\begin{exemple}
Dans une entreprise de $250$ salaries, $90$ travaillent a temps partiel. La proportion vaut
\[
\frac{90}{250}=0{,}36.
\]
Donc $36\pourcent$ des salaries travaillent a temps partiel.
\end{exemple}

\section*{Appliquer un pourcentage}

Prendre $t\pourcent$ d'une quantite revient a multiplier cette quantite par $\frac{t}{100}$.

\begin{exemple}
$18\pourcent$ de $640$ vaut
\[
640\times\frac{18}{100}=115{,}2.
\]
\end{exemple}

\section*{Taux d'evolution}

Lorsqu'une valeur passe d'une valeur initiale $V_i$ a une valeur finale $V_f$, la variation absolue est
\[
V_f - V_i.
\]
Le taux d'evolution est
\[
t=\frac{V_f - V_i}{V_i}.
\]
Il peut s'exprimer sous forme decimale ou en pourcentage.

\begin{exemple}
Un prix passe de $80$ \euro{} a $92$ \euro{}. Le taux d'evolution vaut
\[
t=\frac{92-80}{80}=\frac{12}{80}=0{,}15.
\]
Le prix a augmente de $15\pourcent$.
\end{exemple}

\section*{Coefficient multiplicateur}

Augmenter une valeur de $t\pourcent$ revient a la multiplier par
\[
1+\frac{t}{100}.
\]
Diminuer une valeur de $t\pourcent$ revient a la multiplier par
\[
1-\frac{t}{100}.
\]

\begin{exemple}
Une hausse de $12\pourcent$ correspond au coefficient multiplicateur $1{,}12$.
Une baisse de $30\pourcent$ correspond au coefficient multiplicateur $0{,}70$.
\end{exemple}

\section*{Evolutions successives}

Pour plusieurs evolutions successives, on multiplie les coefficients multiplicateurs.

\begin{exemple}
Un prix augmente de $20\pourcent$ puis diminue de $10\pourcent$. Le coefficient global est
\[
1{,}20\times0{,}90=1{,}08.
\]
L'evolution globale est donc une augmentation de $8\pourcent$.
\end{exemple}

\section*{Evolution reciproque}

L'evolution reciproque permet de revenir a la valeur initiale. Si le coefficient multiplicateur d'une evolution est $C$, alors le coefficient multiplicateur reciproque est
\[
\frac{1}{C}.
\]

\begin{exemple}
Apres une hausse de $25\pourcent$, le coefficient est $1{,}25$. Pour revenir a la valeur initiale, il faut multiplier par
\[
\frac{1}{1{,}25}=0{,}8.
\]
Il faut donc une baisse de $20\pourcent$.
\end{exemple}

\end{document}
